\section*{Limitations}
\label{sec:limitations}

We organize residual limitations by priority for follow-on work:
\begin{itemize}[leftmargin=1.4em,itemsep=1pt,topsep=2pt]
  \item \textbf{P0} --- acceptance-critical: most directly affect the strength of the primary claims under review;
  \item \textbf{P1} --- methodological completeness: improve measurement quality and coverage of the reported experiments;
  \item \textbf{P2} --- longer-horizon generality: domain, jurisdiction, and corpus transfer beyond the current setting.
\end{itemize}

\subsection*{P0 --- Acceptance-critical}

\textbf{P0-1: Statistical power of the human-adjudicated triage claim.}
Because 847 of 1{,}247 labels are retained pre-labels, primary triage claims are restricted to the human-adjudicated test split of 82 cases. On that split the VIGIL--B5 $\Delta$F1 is ${+}0.04$ with 95\% CI $[{-}0.09,{+}0.18]$ (includes zero), so the triage margin over the strongest baseline remains modest and directional rather than statistically established. Expanding the human-adjudicated test set (or reporting a pre-registered larger adjudication) is the highest-priority next measurement.

\textbf{P0-2: Backbone-matched baselines.}
The open-weights re-run (Table~\ref{tab:robustness}) covers only the full VIGIL system, not the baselines. We therefore cannot yet claim that VIGIL's margin over B3--B5 holds independently of the LLM backbone. Backbone-matched comparisons on the same open-weights model remain a required control.

\textbf{P0-3: Cross-stack transfer of the gate.}
The validity-gating operator is defined generally, but all main evidence comes from VIGIL's own pipeline on U.S.\ AI-litigation data. Its transfer to \emph{external} retrieval stacks is specified as Protocol~P1 (Appendix~\ref{app:p1}), which has not been executed in full. A pilot on overruled U.S.\ Supreme Court precedent with gold closures removed stale cases on all three host retrievers but improved nDCG@10 significantly on only one, missing the pre-fixed criterion. The gains reported here should not be read as evidence that the gate improves other retrievers or other domains; completing P1 is a priority generality test.

\subsection*{P1 --- Methodological completeness}

\textbf{P1-1: Extractor recall and edge-error propagation.}
Negative-treatment edge recall (0.77 internal gold; 20.4\% missed-action rate against an authoritative citator, counting missed closures and missed down-weights) is useful but not lossless. A missed action is exactly what allows a stale precedent to leak into a risk decision. VIGIL mitigates this with human-review routing for low-confidence closures, and Table~\ref{tab:edge_sensitivity} quantifies downstream PLRE impact, yet improving extraction recall remains an open methodological problem.

\textbf{P1-2: Benchmark scale and external baselines.}
At the current scale PLRE remains an initial benchmark over five product archetypes and 312 product-mapping cases. We report paired-bootstrap CIs for the full-vs-flat PLRE gap and for a distractor-free subset, but not for every external baseline in Table~\ref{tab:plre_external} or for the intermediate weighting step of Table~\ref{tab:temporal_ablation}. PLRE tables use all 1{,}856 pairs at the case-construction level; bootstrap CIs are computed on the test-split pairs. Since PLRE labels encode temporal validity, part of the gating gain comes from rejecting stale-authority distractors (the gap shrinks when those pairs are removed). PLRE results should therefore be read as evidence for the validity-gated formulation rather than as coverage of all AI-product litigation. The report-quality study (RQ2) uses two blinded raters on 15 cases (mean $\kappa{=}0.68$) without multiple-comparison correction, and is treated as supporting, not primary, evidence. Stronger external legal-RAG / temporal-retrieval baselines and broader archetype coverage would strengthen the benchmark claim.

\textbf{P1-3: Temporal-shift robustness of calibration.}
Proposition~\ref{prop:termination} is a bounded termination-and-escalation property, not a convergence theorem. Its termination clause is close to definitional, and its progress assumption is not verified directly. Calibration controls evidence omission under its calibration distribution (150 cases), not correct legal conclusions, temporal-shift robustness, or freedom from LLM hallucinations~\citep{dahl2024hallucinating}. Because litigation streams are non-stationary, the omission guarantee is weakest under the temporal split (Table~\ref{tab:risk_control_temporal})---precisely where deployment needs it most---so we treat rolling or weighted re-calibration as a deployment requirement rather than an optional refinement.

\subsection*{P2 --- Longer-horizon generality}

\textbf{P2: Coverage, jurisdiction, and domain transfer.}
Cases are counted per docket, so disputes litigated across several related or consolidated dockets are over-represented relative to distinct disputes (App.~\ref{app:inclusion}). The dataset is also skewed toward the largest defendants---the top three account for 41.3\% of cases---limiting representativeness for smaller vendors. VIGIL targets U.S.\ federal CourtListener/RECAP coverage and English-language documents; non-U.S.\ jurisdictions require different docket schemas and citation practices. Modeling partial or jurisdiction-specific treatments, and testing the gate on other superseding corpora (regulations, evolving technical standards), remain longer-horizon open problems; the released SCOTUS testbed (Appendix~\ref{app:p1}) is intended to support the first of these directions.
